\section{Certified Reference Solution}\label{app:ref}

\rev{All accuracy figures of this article are measured against a spectral reference solution computed as follows. The unknowns are the values of $u$ at the Chebyshev--Gauss--Lobatto nodes $x_j=\cos(j\pi/N)$; the two rows $(Du)_0=(Du)_N=0$ of the differentiation matrix $D$ are used to eliminate the wall unknowns (equivalent to a differential--algebraic system with algebraic boundary rows), so that the Neumann condition holds exactly for $t>0$ also when $u_0$ does not satisfy it. Both initial conditions used here are even in $x$, and the equation and the Neumann condition commute with $x\mapsto-x$, so the solver works on the even half ($N/2$ unknowns), which was verified to agree with the full-interval solver to $2\times10^{-13}$. The time integrator is the fifth-order L-stable Radau method (adaptive, analytic Jacobian, rtol $10^{-11}$, atol $10^{-13}$), restarted at every output time so that the returned values are step endpoints. Values at the evaluation points are obtained by barycentric interpolation. The resolution is $N=128$ for $\varepsilon\ge0.05$ and $N=256$ for $\varepsilon=0.01$ ($T\le1$), and $N=\min(2048,\max(256,2\lceil5.12/\varepsilon\rceil))$ for $T>1$.}

\rev{\emph{Certification.} ``Certified'' means demonstrated numerically, not a rigorous a-posteriori bound. (i)~For $T\le1$, $\varepsilon\in\{0.01,0.05,0.1\}$ and both initial conditions, the maximum difference of the default solution at $t\in\{0.1,0.5,1\}$ from the solution with $N=768$ is $\le4\times\nobreak10^{-11}$ (the floor of $\approx3.5\times10^{-11}$ is rounding in the large differentiation matrices); tightening rtol to $10^{-12}$ changes the solution by $\le4\times10^{-13}$ (tested at $N=384$); the BDF and LSODA integrators agree to $\le1.3\times10^{-10}$. (ii)~An independent second-order finite-difference discretization with ghost nodes for the Neumann condition ($N_x=1000,2000,4000,8000$, Radau) converges at the rate $h^2$ and, after Richardson extrapolation, agrees with the spectral solution to $\le5.4\times10^{-9}$ (worst case: the incompatible initial condition at $\varepsilon=0.01$, $t=0.1$) and mostly to $10^{-12}$--$10^{-10}$; since the two discretizations are coded separately, this rules out an error in the sign, scaling or boundary rows of the spectral solver. (iii)~For the interface benchmark B3 ($\varepsilon=0.05$, $T=5$) the default $N=256$ differs from $N=384$, $512$ and $768$ by $2.7\times10^{-12}$, $6.4\times10^{-13}$ and $2\times10^{-11}$ in the maximum norm over $t\in\{0.2,0.5,1,2,3,4,5\}$. (iv)~For $\varepsilon=0.01$ and $T=10$ one needs $N\ge1024$ ($N=256$ is wrong by $10^{-2}$); this case is not used here.}

\rev{\emph{The reference of the first version.} The first version used a semi-implicit finite-difference scheme (first order in time, second order in the interior; $N_x=500$, $\delta t=10^{-4}$) with the wall rows $(1+\lambda)u_0-\lambda u_1$, which are first-order accurate and equivalent to halving the diffusion at the two wall nodes; the ghost-node rows $(1+2\lambda)u_0-2\lambda u_1$ are the consistent ones. Against the certified solution its maximum error on the 200-point grid was $2.7$--$3.7\times10^{-3}$ at $t=0.1$, $1.4$--$2.5\times10^{-3}$ at $t=0.5$ and $5.6\times10^{-4}$--$1.4\times10^{-3}$ at $t=1$ for $\varepsilon\in\{0.01,0.05,0.1\}$ (original initial condition). Re-scoring the stored predictions of the first version against the certified solution changed the $L^2$ errors of the Spectral-PINN by at most $2.2\%$ and those of the vanilla PINN by at most $0.3\%$ (the maximum-norm errors by at most $0.5\%$ and $5.2\%$), so the PINN rows were not materially affected; the classical baseline row was (its error of $8.5\times10^{-5}$ was mostly the error of the reference, and the norm used for that row was inconsistent with the PINN rows: the sum over the $N+1$ Chebyshev nodes was multiplied by the spacing of the finite-difference grid instead of integrating over $[-1,1]$).}

\rev{\emph{Reaction-only comparison.} The closed-form solution of $u'=u-u^3$ at each $x$, $u(x,t)=u_0e^t/\sqrt{1+u_0^2(e^{2t}-1)}$, is used as a reference for ``how much of the solution is explained by the reaction term alone'': its distance to the certified solution is the number given as ``Reaction-only ODE ($\varepsilon=0$)'' in Tables~\ref{tab:forward} and~\ref{tab:orig}.}
